\documentclass[11pt,a4paper]{ctexart}

\usepackage[margin=23mm,headheight=14pt]{geometry}

\usepackage{amsmath,amssymb,graphicx,booktabs,tabularx,longtable,xcolor,fancyhdr,caption,float,placeins}

\usepackage[hidelinks,unicode]{hyperref}

\setmainfont{TeX Gyre Pagella}

\setsansfont{TeX Gyre Heros}

\setmonofont{Consolas}[Scale=.8]

\definecolor{Ink}{HTML}{203746}\definecolor{Teal}{HTML}{197B77}

\ctexset{section={format=\Large\bfseries\color{Ink}}}

\setlength{\parskip}{.4em}\setlength{\emergencystretch}{3em}

\captionsetup{font=small,labelfont=bf}

\pagestyle{fancy}\fancyhf{}\fancyhead[L]{\small MechAgent Lab}\fancyhead[R]{\small AlexanderChan}\fancyfoot[C]{\thepage}

\newcommand{\labfig}[2]{\begin{figure}[htbp]\centering\includegraphics[width=\linewidth]{#1.pdf}\caption{#2}\end{figure}}

\begin{document}\thispagestyle{empty}

{\sffamily\color{Teal}ENGINEERING FIELD NOTES / 2026.11.05}\par\vspace{1em}

{\Huge\bfseries\color{Ink}MechAgent Lab}\par\vspace{.5em}

{\LARGE 力学在环的 Agent 评测}\par\vspace{1em}

从一项力学决策，到原生 Agent、实际求解和可审计产物。

\par AlexanderChan / 2026-11-05\par\vspace{1em}

八道公开题，两个原生客户端，16 次正式尝试。Claude Code 使用 DeepSeek；Codex 使用既有 ChatGPT 登录。只在 S04 运行实际商业求解器，其余题为明确标注的简化力学或工程证据算例。

\tableofcontents\newpage

\section{为什么把力学放进 Agent 评测}

工程仿真中的难点往往不是把命令写出来，而是在物理假设、数值离散和软件状态之间作出一致的选择。一个能够解释子结构原理的模型，未必能够选择可凝聚的区域；一个能读出接触力的模型，也可能把未到终点的中止算例误判为接触失效。因此，这个实验室把评价对象从“回答是否像工程师”改为“决策能否落到可核验的产物”。

\par

MechAgent Lab 从连接器仿真工作流中抽取八个公开任务：量纲与平衡、网格分配、小板预抬、子结构分区、有界公差、接口复用、偏位偏角边界和批次恢复。题目使用原创公开杆件及简化力学数据，不上传客户模型、合同或工业连接器结果。

\par

本轮新增两个真实原生执行循环：Claude Code 接入既有 DeepSeek 路线，Codex CLI 使用既有 ChatGPT 登录。CLI 是 Command-Line Interface，即命令行界面；CCS 在本文中指本机 CC Switch 配置管理工具，不是模型或 Agent。此前受控 mini/Kimi 与 Pi/DeepSeek 的轨迹仍可阅读，但它们不是本轮相同矩阵上的重跑成绩。框架、模型和服务路线是三个不同的变量，不能因为名字叫 Claude Code 就把成绩归给 Claude 模型。

\par

\labfig{architecture}{模型、原生框架、受控工具和主机判分分层；mini 与 Pi 为历史路线。}

\section{先约定证据，再讨论成绩}

每个 Agent 只能通过三个 MCP 工具读取题目、执行实验和提交结构化答案。MCP 的全称为 Model Context Protocol，中文为模型上下文协议；这里它是工具调用的接口，不是力学求解器本身。工具参数中没有任意脚本、文件路径或终端命令。参考解与判分器由主机持有，不通过工具开放。

\par

原生表示各客户端仍自行完成推理、工具选择和上下文管理，而不是把其模型装进自制循环。Claude Code 关闭无关工具、插件及全局配置发现；Codex 关闭终端、浏览器和其他无关能力，只使用显式配置的力学 MCP。这里声称的是配置约束和可观察工具审计，不是操作系统级的恶意程序隔离。

\par

代码、题面和判分规则在正式调用前冻结并保存 SHA256。每个客户端对每题只保留一次正式尝试；模型失败不追加调参重试。两次 Claude 工具未加载和七次 Codex 握手失败的启动记录完整保留并单列排除，已经有效的 Codex 第一题及全部 Claude 正式记录不重跑。连接修正只涉及接入层，没有修改力学题面或答案。

\par

预算也必须公开：每题最多 24 次受控工具调用和 360 秒墙钟时间，子结构题最多一次候选实际求解。Claude 另有限定的原生轮次及模型请求输出上限；Codex 未声称有相同的输出 token 硬上限。不同默认系统指令和预算意味着，本轮可以描述发生了什么，但不能据此推出纯粹的模型或框架因果排名。

\par

\section{八道题如何检查力学而不是措辞}

S01 采用串联轴向杆，要求先把 GPa 转成 MPa，再独立核对端位移、根部反力和应变能。有限元法的英文全称为 Finite Element Method，简称 FEM。评分不仅看最终数字，还检查数字是否对应真实执行过的工具产物。

\par

S02 使用截面积随位置变化的杆件。Agent 需要选满足端位移误差限值的最少单元数，而不是无条件使用最细网格。独立参考来自连续体柔度积分，离散模型则使用单元中点面积。该设计将准确性和计算资源选择联系起来。

\par

S03 将六块小板建成六个独立控制量，每块影响两对单边接触。抬升在名义值上叠加，目标为所有法向力不低于 0.4 N，同时使幅值平方和最小。单边弹簧在分离时返回零法向力，因此不能用摩擦力来补足接触阈值。该题是可独立核验的简化接触题，而不是三维摩擦接触仿真。

\par

S04 才是本轮真实商业求解器题。原创三维杆件一半为线弹性材料，另一半为 J2 弹塑性材料。Agent 在候选截面中选择可凝聚区域，保留根部和接口的三向平移，并实际生成子结构、装配求解、提取 ODB 数据。ODB 的全称为 Output Database，中文为结果数据库。

\par

S05 生成 72 维公差样本，包括 42 个 Z 向形状变量、6 个小板预抬和 24 个公端子 X/Y 局部平移。MC 是 Monte Carlo（蒙特卡洛），LHS 是 Latin Hypercube Sampling（拉丁超立方抽样），Sobol 与 Halton 是低差异序列。四种方法先生成单位超立方体样本，再用截断标准正态的逆累积分布映射到 [-1,1]，而不是把超界正态数裁剪到端点。

\par

物理映射不能混淆：形状变化为 0.15a mm，小板预抬为 0.3348+0.15a mm，公端子的刚体局部平移为 0.1a mm。局部制造偏移不移动子结构，也不改变端子尺寸；整套公端相对母端的装配偏位是另一种姿态变量。

\par

S06 同时检查子结构库身份和 Tie 接口覆盖。CORE 指仍以完整有限元描述的区域，SUB 指凝聚为子结构的区域；二者不是特定软件实体名称。只有拓扑、材料和保留节点坐标都一致的库才能复用。Tie 位置容差只扩大几何搜索范围，不能修复错误材料或错误网格；关闭节点调整，才能避免把制造偏移吸回原位置。

\par

S07 沿指定的偏位偏角耦合路径寻找首个失效边界。每次实验计算 12 对法向力，取最小值与 0.4 N 比较；最终需要实测安全端与失效端两个产物，而不是只给一个看似收敛的角度。这个题目使用刚体杠杆与单边弹簧，不把其数值当成工业连接器极限。

\par

S08 聚焦容易被忽视的工程可信性。只有求解成功、达到最终位移且 12 个通道完整的 case 才进入排名；case 指一组输入及其对应结果。对每个有效 case 先取绝对法向力的最小值，再按该值升序排序。物理接触失效的已完成算例不重算；活动任务等待，缺通道只重提取，中断任务重算，负体积输入先修网格。

\par

\labfig{tasks}{八题的决策对象及证据层级，只有 S04 调用真实 Abaqus。}

\[u=F\sum_i\frac{L_i}{E_iA_i},\quad R=-F,\quad W=\tfrac12Fu\]

\[a=\Phi^{-1}\!\left[\Phi(-1)+u\bigl(\Phi(1)-\Phi(-1)\bigr)\right],\quad u\in(0,1)\]

\labfig{sampling}{四种有界抽样的代表变量示意；16 个样本不用于比较统计效率。}

\[g(\lambda)=\min_{j=1,\ldots,12}N_j(\lambda)-0.4\;\mathrm{N}\]

\labfig{boundary}{S07 简化力学路径与实际 Agent 实验点；曲线不是工业连接器预测。}

\section{真正运行子结构：少保留哪些自由度}

三维公开杆件尺寸为 12×2×2 mm，共 117 个节点和 48 个 C3D8 单元。可线性凝聚的范围不能跨入可能塑性的区域。消去内部自由度后，界面刚度由 Schur 补得到；DOF 是 Degree of Freedom，即自由度。最大安全截面为 6 mm，消去 135 个内部平移自由度。

\par

完整模型先独立求解，Agent 不能读取该参考。它提交分区后，主机工具通过真实 Abaqus/MCP 生成 SIM 子结构库、进行降阶装配并求解，再将存储帧反力曲线与完整模型比较。微型算例用于证明执行和等价性，不用于宣称工业级加速比。图中曲线来自本轮成功的 Claude Code/DeepSeek 实际提取，不是理论曲线替代软件计算。

\par

Codex 同样选出 6 mm 截面和 135 个可消去自由度，但实际工具在 240 秒内未返回产物；再次调用被单候选预算拒绝。它明确提交空产物编号而没有捏造求解成功。其独立后台求解随后完成，但没有在尝试结束前返回给模型，也没有完成该次曲线提取核验。原始验收零分保留，单列为工具超时／证据未完成；后台迟到的成功不能补算 Agent 通过，也不替它画一条已验证曲线。

\par

\[K_c=K_{bb}-K_{bi}K_{ii}^{-1}K_{ib}\]

\labfig{solid}{本轮原生 Agent 的实际完整／降阶反力曲线与公开分区。}

\section{怎样阅读结果和失败}

结果矩阵的每个格子对应一个已冻结的“题目×原生客户端”组合。通过必须同时满足独立数值检查和产物引用；预算耗尽、没有提交、服务错误和接入失败分别记录。不能因为工具文本写着成功，就把未完成求解当成证据。

\par

耗时包含客户端启动、模型等待和实际工具执行，S04 还包含真实商业求解器阶段。它不是纯模型推理速度，也不是求解器加速比。费用字段同样区分 SDK 估算与真实账单；第三方路线的 SDK 美元估算不一定对应实际收费，未知账单保持为空。

\par

失败轨迹比单一通过率更能说明工程问题：正确推导但未提交是闭环失败，完整数据映射错误是物理语义失败，服务限流是环境问题，MCP 未加载是接入问题。本轮不通过追加尝试掩盖这些差异。

\par

Claude 的 S07 找到了安全端 1.223 和失效端 1.224，却以含 safe/fail 键的对象提交产物，而冻结判分器按数组读取。题面没有显式限定这一字段为数组，因此原始零分保留，并将该格标为协议歧义，不解释为接触力边界判断错误。它没有获得追加尝试或事后满分。

\par

正式尝试中，通过 12 次，数值／语义未通过 1 次，预算耗尽 1 次，环境或未完成 1 次，协议歧义 1 次。另保留 9 次接入排除。

\labfig{outcomes}{逐题状态及墙钟耗时；不同模型、路线和预算不能作因果排名。}

\begin{longtable}{llrrl}\toprule Task & Client & s & Tools & Outcome\\\midrule\endhead

S01 & CC & 17.5 & 0 & Setup excluded \\

S01 & CC & 12.7 & 0 & Setup excluded \\

S01 & CC & 11.5 & 4 & Pass \\

S02 & CC & 18.0 & 5 & Pass \\

S03 & CC & 17.2 & 7 & Budget \\

S04 & CC & 64.4 & 3 & Pass \\

S05 & CC & 19.7 & 6 & Fail \\

S06 & CC & 9.6 & 3 & Pass \\

S07 & CC & 16.9 & 18 & Contract excluded \\

S08 & CC & 10.4 & 4 & Pass \\

S01 & Codex & 155.7 & 3 & Pass \\

S02 & Codex & 6.6 & 0 & Setup excluded \\

S02 & Codex & 39.8 & 3 & Pass \\

S03 & Codex & 2.0 & 0 & Setup excluded \\

S03 & Codex & 56.7 & 4 & Pass \\

S04 & Codex & 2.0 & 0 & Setup excluded \\

S04 & Codex & 294.4 & 5 & Incomplete \\

S05 & Codex & 2.0 & 0 & Setup excluded \\

S05 & Codex & 60.7 & 6 & Pass \\

S06 & Codex & 2.0 & 0 & Setup excluded \\

S06 & Codex & 30.6 & 3 & Pass \\

S07 & Codex & 2.0 & 0 & Setup excluded \\

S07 & Codex & 48.5 & 4 & Pass \\

S08 & Codex & 3.0 & 0 & Setup excluded \\

S08 & Codex & 34.8 & 3 & Pass \\

\bottomrule\end{longtable}

\section{复现：从一条工具调用追到一份产物}

下载包包含题卡、冻结源码、四份连续注册记录、所有保留尝试、逐轮输入输出附录、独立判分、原创公开 INP、实际 MCP 任务标识、绘图源数据以及 Tavotto 导出回执。所有图均在实际 Tavotto 中执行、注册并导出为矢量 PDF 和 PNG。正文由 LaTeX 编译，中英文各保留一份。

\par

复现前先运行离线测试和本地假服务的工具加载检查，再配置自己的 API 或原生登录、MCP 环境及商业软件许可证。不要把密钥写入题目或上传 CCS 数据库。注册完成后不修改题面、工具和判分器；若需要改变协议，应另建记录并保留旧证据。

\par

这个基准的下一步不是把八道题都包装成商业软件任务，而是逐步增加真正需要三维接触、网格质量和路径依赖的公开算例，并建立重复、交叉配置和费用受控的研究设计。在此之前，清楚说明证据来自哪里，比给出一个漂亮排行榜更重要。

\par

\section{资料与实现依据}

\url{https://learn.chatgpt.com/docs/non-interactive-mode}\par

\url{https://code.claude.com/docs/en/headless}\par

\url{https://alexchan114.github.io/post/2026/11/04/mechanics-benchmark-v02/}\par

\end{document}